% ══════════════════════════════════════════════════════════════════════════════
%  MINI-LAB A — THE VISCOSITY RACE                                     (Lesson 1)
%  A short class investigation inside Lesson 1, lettered so that Experiments
%  1–3 keep the numbers the decks and the lesson booklet already use.
% ══════════════════════════════════════════════════════════════════════════════
\renewcommand{\thesection}{A}\renewcommand{\symid}{0}
\renewcommand{\expword}{Mini-lab}\renewcommand{\expcaps}{MINI-LAB}
\experimentdiv{The Viscosity Race}{Planning a fair test \textbullet{} Timing with a stopwatch \textbullet{} Ranking liquids by how easily they flow}{1}\label{sec:A}
\lessoncodes{VC2S8U05 \textbullet{} VC2S8I01--I05}

\vocabhead
\begin{ctab}
\begin{tabular}{|SW{4.2cm}|W{11.6cm}|}
  \hline\hd & \thead{Word} & \thead{Definition (copy from the slide)} \\ \hline
  & \scaf{\textbf{Viscosity}} & \scaf{How thick a liquid is: how much it resists flowing. Honey has a high viscosity.} \\ \hline
  & \textbf{Fair test} & \\ \hline
  & \textbf{Variable} & \\ \hline
  & \textbf{Trial} & \\ \hline
  & \textbf{Average} & \\ \hline
\end{tabular}
\end{ctab}

\begin{lipanel}
\inquiryli{Plan a fair test, and time each liquid with a stopwatch.}
\understandli{Compare how easily liquids flow, and rank them by viscosity.}
\sctitle
\begin{itemize}[leftmargin=1.4em, itemsep=0pt, topsep=1pt]
  \item I can name the variables I keep the same to make the race a fair test.
  \item I can time each liquid three times and calculate an average.
  \item I can rank liquids from lowest to highest viscosity, using my results as evidence.
\end{itemize}
\end{lipanel}

\begin{skillsfocus}
\focus{QP1}{Predict the order the liquids will finish in, and say why.}
\focus{PC2}{Decide what to change, what to measure and what to keep the same.}
\focus{PC6}{Time three trials for each liquid and record them straight away.}
\focus{PC8}{Start and stop the stopwatch at the start and finish lines.}
\focus{PA1}{Record your times in a table with units in the headings.}
\focus{PA5}{Describe how race time relates to viscosity.}
\end{skillsfocus}

% ─── TECHNIQUE ────────────────────────────────────────────────────────────────
\exppart{Technique spotlight}{Blue}

\begin{technique}{Making a race a fair test}{PC2}
A test is \textbf{fair} when only one thing is different between the runs, so any difference in the results must come from that one thing.
\begin{itemize}[leftmargin=1.4em, itemsep=1pt, topsep=3pt]
  \item The \textbf{independent variable} is what you change on purpose: here, \emph{which liquid} races.
  \item The \textbf{dependent variable} is what you measure: here, \emph{the time to reach the finish line}.
  \item The \textbf{controlled variables} are everything you keep the same: the amount of liquid, the slope of the tile, the distance between the lines, the room temperature.
\end{itemize}
Run each liquid more than once. Each run is a \textbf{trial}; the average of three trials is more reliable than one.
\end{technique}

\begin{workedex}{Variables in a toy-car race}
Three toy cars are rolled down the same ramp to find which is fastest. Identify the variables.

\smallskip
\textbf{Independent:} which car is used. \quad \textbf{Dependent:} the time to reach the bottom.\\
\textbf{Controlled:} the same ramp, the same slope, the same starting line, and letting go without pushing.
\end{workedex}

\begin{wedo}{Variables in the viscosity race}
Complete the table for the viscosity race.
\begin{ctab}
\begin{tabular}{|SW{4.0cm}|W{10.6cm}|}
  \hline\hd & \thead{Variable} & \thead{In the viscosity race} \\ \hline
  & Independent & \\ \hline
  & Dependent & \\ \hline
  & Controlled (1) & \\ \hline
  & Controlled (2) & \\ \hline
\end{tabular}
\end{ctab}
\end{wedo}

\begin{technique}{Timing with a stopwatch}{PC8}
Share the jobs: one person \textbf{releases} the liquid, one \textbf{times}, one \textbf{records}. The timer starts the stopwatch the moment the front of the liquid crosses the start line, and stops it the moment it crosses the finish line. Record the time to one decimal place, in seconds (s). Reset the stopwatch before the next trial.
\end{technique}

\begin{youdobox}[Averages]
\task Calculate the average time for each liquid. Round to one decimal place. \skr{PA3}

\begin{ctab}
\begin{tabular}{|SW{3.4cm}|Z{2.3cm}|Z{2.3cm}|Z{2.3cm}|Z{3.6cm}|}
  \hline\hd & \thead{Liquid} & \thead{Trial 1 (s)} & \thead{Trial 2 (s)} & \thead{Trial 3 (s)} & \thead{Average (s)} \\ \hline
  & \scaf{e.g. milk} & \scaf{2.1} & \scaf{2.4} & \scaf{2.2} & \scaf{$6.7 \div 3 \approx 2.2$} \\ \hline
  & Shampoo & 18.2 & 17.6 & 18.4 & \\ \hline
  & Tomato sauce & 31.0 & 29.4 & 30.2 & \\ \hline
\end{tabular}
\end{ctab}
\end{youdobox}

% ─── PLAN AND SAFETY / METHOD ────────────────────────────────────────────────
\exppart{Plan, safety and method}{Teal}

\textbf{Aim:} To compare the viscosity of four liquids by timing how long each takes to run down a sloping tile.

\needspace{10\baselineskip}
\task Predict the order the four liquids will finish in, from fastest to slowest, and explain your prediction. The liquids are water, cooking oil, dishwashing liquid and honey. \skr{QP1}
\answerbox{5cm}

\subsubsection*{Equipment}
\begin{multicols}{2}
\begin{itemize}[leftmargin=1.2em, itemsep=0pt, topsep=0pt]
  \item smooth white tile (or plastic tray)
  \item two textbooks, to prop the tile up
  \item ruler and whiteboard marker
  \item four 5\,mL measuring spoons (one per liquid)
  \item water, cooking oil, dishwashing liquid, honey
  \item stopwatch
  \item paper towel
\end{itemize}
\end{multicols}

\begin{center}
\begin{tikzpicture}[scale=0.9]
  \draw[eq] (-0.2,0) -- (9.2,0);
  \draw[eq] (0,0) rectangle (2.0,0.55) (0.1,0.55) rectangle (1.9,1.1);
  \draw[eq] (0.3,1.25) -- (8.6,0.05) -- (8.5,-0.05) -- (0.25,1.1) -- cycle;
  \draw[Ink, line width=1.2pt] (1.95,1.03) -- (2.1,1.2);
  \draw[Ink, line width=1.2pt] (6.6,0.37) -- (6.75,0.55);
  \fill[black!70] (1.6,1.08) ellipse (0.16 and 0.07);
  \eqlab{(2.05,1.15)}{(2.3,2.2)}{start line}
  \eqlab{(6.7,0.5)}{(7.0,1.6)}{finish line, 20\,cm down the tile}
  \eqlabl{(1.0,0.85)}{(-0.3,2.0)}{textbooks}
  \eqlab{(4.5,0.65)}{(4.6,1.9)}{tile}
\end{tikzpicture}
\expfigcap{Scientific diagram of the viscosity race.}
\end{center}

\begin{infobox}{Safety}
Spilt oil, honey and dishwashing liquid make the bench and floor slippery: wipe spills up straight away. Never taste anything in the lab. Wash your hands at the end.
\end{infobox}

\subsubsection*{Method}
\begin{enumerate}[itemsep=2pt]
  \item Prop one end of the tile on two textbooks. Draw a start line near the top and a finish line 20\,cm further down.
  \item Fill a measuring spoon level with the first liquid.
  \item Pour the liquid just above the start line. The timer starts the stopwatch as the front of the liquid crosses the start line.
  \item Stop the stopwatch as the front of the liquid crosses the finish line. Record the time.
  \item Wipe the tile clean and dry. Repeat for two more trials.
  \item Repeat steps 2--5 for each of the other liquids, keeping the tile at the same slope.
\end{enumerate}

% ─── RESULTS ─────────────────────────────────────────────────────────────────
\exppart{Results}{Amber}

\subsubsection*{Time for each liquid to run 20\,cm down the tile \skr{PA1}}
\begin{ctab}
\begin{tabular}{|SW{3.4cm}|Z{2.1cm}|Z{2.1cm}|Z{2.1cm}|Z{2.6cm}|Z{2.2cm}|}
  \hline\hd & \thead{Liquid} & \thead{Trial 1 (s)} & \thead{Trial 2 (s)} & \thead{Trial 3 (s)} & \thead{Average (s)} & \thead{Rank (1 = fastest)} \\ \hline
  & \scaf{e.g. milk} & \scaf{2.1} & \scaf{2.4} & \scaf{2.2} & \scaf{2.2} & \scaf{---} \\ \hline
  & Water & & & & & \\ \hline
  & Cooking oil & & & & & \\ \hline
  & Dishwashing liquid & & & & & \\ \hline
  & Honey & & & & & \\ \hline
\end{tabular}
\end{ctab}

\needspace{10\baselineskip}
\task Write down anything you observed while the liquids ran, other than the times (for example, the shape of the liquid or how it spread).
\answerbox{5cm}

% ─── ANALYSIS ────────────────────────────────────────────────────────────────
\exppart{Analysis and evaluation}{Coral}

\task List the four liquids from lowest viscosity to highest viscosity. \skr{PA5}
\answerbox{1.2cm}

\needspace{10\baselineskip}
\task Describe the relationship between a liquid's race time and its viscosity. Use two liquids from your results as evidence. \skr{PA5}

\scaf{The longer a liquid took to reach the finish line, the \dots. For example, \dots\ took \dots\,s but \dots\ took \dots\,s.}
\answerbox{5cm}

\needspace{10\baselineskip}
\task Was your prediction correct? Use your results to explain. \skr{EV1}
\answerbox{5cm}

\needspace{10\baselineskip}
\task Explain why each liquid was timed three times instead of once. \skr{PC6}
\answerbox{5cm}

\needspace{14\baselineskip}
\task A group used a full spoon of honey but only half a spoon of water. \skr{PC2}
\pt{a}Which controlled variable did they change?
\answerbox{1cm}
\pt{b}Explain why this makes their race unfair.
\answerbox{5cm}

\begin{yr8box}
\begin{keepext}
\exttask Honey flows much more easily when it is warm. Predict how the honey's race time would change if the honey and the tile were warmed to 40\,\textdegree C, and design a fair test to check your prediction.
\answerbox{5cm}
\end{keepext}
\begin{keepext}
\exttask Toothpaste does not flow down the tile at all unless it is pushed, but it is not a solid either. Suggest how you could decide whether toothpaste is a solid or a liquid.
\answerbox{5cm}
\end{keepext}
\end{yr8box}

% ─── SKILL CHECK ─────────────────────────────────────────────────────────────
\exppart{Skill check}{NavyBlue}

\begin{skillcheck}
\scrow{QP1}{I can predict an order and give a reason for my prediction.}
\scrow{PC2}{I can name the variables I change, measure and keep the same in a fair test.}
\scrow{PC6}{I timed three trials for each liquid and recorded them straight away.}
\scrow{PC8}{I started and stopped the stopwatch at the lines.}
\scrow{PA1}{I recorded my results in a table with units in the headings.}
\scrow{PA5}{I can describe how race time relates to viscosity, using data.}
\end{skillcheck}

\task Which skill will you focus on in Experiment 1? Explain what you will do differently.
\answerbox{5cm}

% back to numbered experiments
\setcounter{section}{0}\renewcommand{\thesection}{\arabic{section}}\renewcommand{\symid}{\thesection}
\renewcommand{\expword}{Experiment}\renewcommand{\expcaps}{EXPERIMENT}
